\section{Preserving a model before extending its accounts}
\label{sec:company-module-language}

The scenario family now keeps each company's instructions and consequences
together. Exchanging those results with another implementation raises two
further questions. Can it identify the same input despite a different text
encoding? Can it add a new accounting method without creating another cash
balance or losing the company's tax history? The first is a representation
question. The second requires a richer account of company state.

\paragraph{Different bytes can express the same instruction.}
A company name containing an accented character may be written directly in
UTF-8 or through a JSON escape. Both decode to the same sequence of characters,
but their original files differ. Keeping the exact source file and identifying
the decoded model therefore require separate identities.

The legacy import retains the original bytes and their hash, the declared local
decoder and profile, and the existing company draft's decoded identity. A
separate portable content record binds the canonical decoded value. Equivalent
escaping or whitespace can share that content record while retaining distinct
original imports. Different character sequences are not normalized into one
name. The exchange itself changes when its original bytes change.

A field extracted from a saved result is not passed off as an original file.
Its import identifies the parent and exact JSON pointer, and validation checks
that the value is present there. Family-member records expose names, paths,
results, contexts and annual outcomes; expression records connect the program,
policy, selected parent, generated child and declared dependencies. Recovery
records retain the same family and method with the completed member prefix.
Missing dependencies and wrong reference kinds prevent import. This is an
explicit reference to legacy drafts, rather than a conversion of every field
into the full portable numeric and type language.

Identity checks alone cannot establish that cash was calculated correctly.
One adverse example changes a cash amount and recomputes the affected hashes,
including the saved recovery result. The resulting packet is internally
consistent and passes content binding. Repeating the calculation with the
trusted local engine rejects it. The two verdicts remain distinct: a complete
dependency graph preserves what was supplied, while semantic replay checks
what the selected method computes. Neither supplies an actual source's
publication date, proves historical execution or establishes physical storage
durability. Executable methods are selected by the caller; packet content
cannot load code.

\paragraph{A module uses a company's existing book.}
Consider the retained return-right example. Its legal book already contains
cash 100, a recovery asset 16 and a refund liability 40. A return module must
refer to those balances. Adding its opening cash or capital to the company
again would create resources that no event supplied.

The finite invocation therefore binds the original entity, legal book,
branch, source context, information view, method and complete opening state.
The trusted return method produces a proposed state transition. It uses the
named book's cash, preserves other legal books and exposes the journal and
gross cash effects. The checked example pays 48 and ends with cash 52. No
commit is implied by the proposal. Changing the owner, method, source, state,
event or account map prevents reuse; missing and false findings remain
distinct from a successful calculation.

The retained owner is \code{Synthetic-return-expected}, with settlement on
30 June 2024. It is not Northstar. Northstar's current language has no
return-right memory, refund/recovery accounts or return-settlement operation.
Its homogeneous inventory and immediate-deductible tax method also cannot
represent this case unchanged. The successful finite proposal identifies a
usable method boundary; native company support needs the state design below.
The separate sector-company method remains report-only and cannot be turned
into a posting operation by relabelling its result.

\paragraph{Returned goods need their own carrying value.}
The original return contract sold ten units at cost 12, with a refund of 20
per returned unit. Its opening estimates recognize a refund liability 40 and
recovery asset 16. These amounts depend on the supplied expected returns,
handling and deterioration. The historical accounting source requires the
refund liability and recovery asset to be distinguished, with the latter
based on former carrying amount less expected recovery costs and loss of
value \citep[IFRS~15.B21--B27]{eu20231803}.

At final settlement, two units return. Their deterioration is 2 per unit,
handling costs 4 per unit and the supplied net realizable value of the lot is
20. Handling is assumed to be selling administration; costs needed to bring
inventory to its location and condition would require a different treatment
\citep[IAS~2.15--16]{eu20231803}. The recovered lot is therefore carried at
$\min\{2(12-2),20\}=20$, and cash payments are 40 for refunds and 8 for handling.

The settlement debits the refund liability 40 and inventory 20, credits the
recovery asset 16 and cash 48, credits cost of sales 4 and debits selling
expense 8. The journal balances and records a book loss of 4. A later sale
of the whole recovered lot for 40 gives revenue 40 and cost of sales 20.
The lot retains its original contract, quantity, carrying amount, tax base
and date so this second event can consume the right stock once.

Simply adding two units to Northstar's homogeneous inventory would value them
at its existing unit cost, which differs from the recovered lot's cost of 10.
The extended book instead reconciles the original homogeneous lot and each
recovered lot separately. Let $q_0$ and $c_0$ denote the homogeneous units and
unit cost, and let $V_\ell$ denote each recovered lot's carrying amount, all
for the same legal book and currency. The proposed inventory control is
\begin{equation}
 I=q_0c_0+\sum_{\ell}V_\ell.
 \label{eq:module-return-inventory}
\end{equation}
The sum preserves different carrying values without silently selecting a
cost-flow method. With no recovered lots, equation~\eqref{eq:module-return-inventory}
reduces to the existing inventory rule. Whole-lot sale is the initial subset;
partial sales and changes of estimates require further specified methods.

\paragraph{Tax history changes the cash needed at the close.}
The same return appears in the joined company case of
Section~\ref{sec:company-composition}. Its legal parent P supplies a complete
trace through settlement, later sale and two tax closes. P's opening, dates
and tax premises are retained; they are not new facts about Northstar.
The case assumes that the original sale and cost were taxed and deducted,
refunds and handling become deductible on settlement, and recovered inventory
enters tax stock at its carrying amount. The supplied tax rate is 25 percent.
The required enactment and recoverability findings remain explicit. Tax bases,
recognition conditions and recovery manner matter independently of book profit
\citep[IAS~12.7--10,24--25,46--51]{eu20231803}.

To see the difference, let $F$ and $R$ be the opening refund liability and
recovery asset, $n$ actual returned units, $u$ refund per unit, $h$ selling
handling per unit and $V$ recovered stock. Under these particular tax premises,
the settlement's book profit $B$ and taxable profit $T$ are
\begin{align}
 B&=F-nu-R+V-nh, & T&=-nu-nh+V,\\
 T-B&=R-F. &&
 \label{eq:module-return-tax}
\end{align}
The tax deduction still includes the refund whose book expense was recognized
before the opening; the book reversal of recovery value has its own earlier
tax difference. Substituting $F=40$, $R=16$, $nu=40$, $nh=8$ and $V=20$ gives
$B=-4$, $T=-28$ and the one-time adjustment $T-B=-24$.

The opening liability's zero tax base gives a deferred tax asset of 10, while
the recovery asset's zero tax base gives a deferred tax liability of 4. Their
reversal adds deferred tax expense 6. P's first-half pretax profit 56 therefore
produces taxable income 32, current tax 8 and after-tax profit 42. Annual
pretax profit 66 gives taxable income 42 and total current tax 10.5; the
second close recognizes only the remaining 2.5. Applying the $-24$ adjustment
afresh to each interval would deduct it twice. The tax history binds the
settlement event and subtracts tax already recognized from the cumulative
calculation.

Multiplying first-half book profit by 25 percent would instead give current
tax 14 and omit the separate deferred expense 6. A coincidentally equal total
expense would not make the current cash payment correct. A complete native
tax close consequently needs tax bases, timing, rate, origin, recovery support
and the full temporary-difference inventory. Without them the tax-dependent
annual report remains unresolved, while earlier valid journals survive as a
dated prefix. The example establishes no actual tax statute.

\paragraph{The statement bridge must retain opening obligations.}
P's annual profit is 49.5, but operating cash is 101.5. Depreciation 20 and
deferred tax expense 6 are noncash charges. Inventory falls by 60, the recovery
asset falls by 16 and the refund liability falls by 40. Another payment is
easy to overlook: the matured opening debt includes 10 of interest expensed
before this year. The selected historical interest policy classifies that
payment as operating, although current finance expense is zero
\citep[IAS~7.31--33]{eu20231803}. The exact indirect bridge is
\begin{equation}
 49.5+20+6+60+16-40-10=101.5.
 \label{eq:module-return-cash}
\end{equation}
Omitting the last term overstates operating cash by 10. Subtracting the whole
debt payment instead would move 90 of financing principal into operating
cash. Equation~\eqref{eq:module-return-cash} agrees with the separate gross
cash events. Opening cash 500, operating cash 101.5, investing cash $-20$ and
financing cash $-90$ yield closing cash 491.5. Opening equity 622 plus profit
49.5 yields 671.5.

These are P's legal-book amounts. Its internal receipt from S remains a legal
cash inflow and is eliminated only in the group report. The same distinction
prevents an apparently balanced module from double-counting internal flows.
Deferred tax is added back once in the indirect bridge; separately adjusting
DTA and DTL again would repeat the same effect.

\paragraph{What the native extension must carry.}
The richer profile needs closed account and component schemas, explicit
return/lot authoring, strict eligibility findings, one chronological shared
journal executor, tax-origin history and extended report mappings. Every
temporary account must enter profit and the closing transfer to retained
earnings. A complete year carries remaining lots, rights, tax bases, event
identities and source context forward. It resets period income and cash-flow
totals without deleting the history needed to prevent a second settlement.
A partial year supplies no complete next-year opening.

The existing company profile remains unchanged. Selecting the richer profile
is explicit; it cannot redefine old record fields or silently replace a
method. The design and exact legal-book trace specify the required repair.
Section~\ref{sec:company-return-reference} now supplies a closed component
authoring profile and a finite reference for return, tax, legal reports and
carry. A complete composed-company language and production integration remain
unfinished. Broader
corporate, group, bank and insurance methods retain their separate obligations.

The supporting reference records are under \path{data/accounting-master/}:
\begin{quote}\small
\path{company-portable-authoring-20261023}\\
\path{company-module-boundary-20261023}\\
\path{company-native-return-design-20261023}
\end{quote}
They preserve original bytes, a finite shared-book proposal and the complete
retained legal-book comparison. Independent implementation, domain and human
acceptance remain pending.

\section{From a return declaration to the next annual opening}
\label{sec:company-return-reference}

P's recovered stock gives the language a concrete implementation problem.
The return contract, the legal book and the tax calculation each remember
something different about the same two units. The contract records that the
customer's rights have ended. The stock record keeps the recovered carrying
amount until sale. The tax calculation remembers an adjustment arising from
expense recognized before the year began. A company program must preserve
all three without posting the opening balances again.

The finite reference now connects these records through the complete retained
2025 P trace. It implements return settlement, whole-lot sale, P's selected
current and deferred tax method, legal statements and annual carry. Other
corporate events continue to use the checked retained methods. This division
makes the implementation comparison precise: the new return and tax
calculations can be compared with the earlier company case, while the
surrounding methods supply the same corporate transactions as before.

\paragraph{An author declares facts before choosing an event.}
The versioned \code{ReturnModule}\allowbreak\code{AuthoringDraft/1} describes a component of
one legal book. It identifies the source packet and opening state, entity,
currency, period and account roles. Its contract memory includes sold and
expected-return units, unit refund and cost, estimated recovery costs, and
outstanding refund and recovery balances. The recovered-stock declarations
keep quantities, carrying amounts and tax bases separately. Tax declarations
state the selected timing method, rate and required findings.

These declarations allocate existing accounts. For example, a refund balance
of 40 must agree with the liability in P's opening book; it is not an
instruction to create another liability. The deferred asset and liability
must likewise agree with their declared opening allocations. The remaining
company balances come from the shared opening, rather than from a second
component wallet.

Typed records and text declarations contain the same fields. The text form
uses named declarations for the component, lots, tax premises and dated
settlement or sale. It contains no executable Python, external include or
opaque JSON body. Both forms preserve canonical exact rational amounts, strict
Boolean findings and source references. Rendering a typed declaration as text
and reading it again must reproduce the same value. This is a matched
authoring comparison with the typed-library baseline, not evidence that text
is preferable for every user.

Unknown values are explicit. Suppose the tax rate has not been established
but the author also supplies a refund liability of 41 while two expected
returns at 20 require 40. The contradiction is already decidable and the
record is invalid. The unknown tax rate cannot conceal it. If the amounts
agree but a required recovery finding is false, the selected method is
unsupported. If the required finding is unknown, the affected calculation
remains unresolved. Arithmetic checks inspect their own dependencies rather
than replacing unknowns with zero or postponing every known check.

Source correspondence is another question. A syntactically valid source
pointer may name the wrong disclosure. The retained example therefore checks
the entire declared projection against P's original packet, including owner,
dates, amounts and opening identity. Changing P to Northstar fails that
correspondence even though both names are valid text. A correctly bound local
packet still needs separate evidence for external authenticity and any
historical admission claim.

\paragraph{One event changes one shared book.}
The reference state combines P's account balances with return memory,
recovered lots, dated tax adjustments, classified gross cash and lifetime
event identities. Every account belongs to one financial-statement class;
every temporary account enters profit. A proposed return first checks its
component, quantities, findings and date. It then derives the settlement
journal already shown in Section~\ref{sec:company-module-language}, creates
the separate recovered lot and closes the return rights. The surrounding
inventory allocation stays intact.

The proposal identifies the complete preceding state, the exact event and
the selected method bundle. Applying it repeats the calculation with the
trusted local method and compares the journal, gross cash and complete
after-state. A changed cash amount cannot become valid merely by supplying a
new hash. A stale proposal, repeated settlement or sale of an already disposed
lot is refused. This is an in-process reference for checking a transition;
durable publication and concurrent writers require the separate storage design.

Cash sufficiency is checked after deriving the complete proposed journal.
If P cannot pay the refund and handling costs, the event returns a funding
failure and leaves the preceding book intact. It does not create financing.
Missing tax evidence has a different consequence: the valid non-tax events
can remain as an exact dated prefix, but the tax-dependent annual statement
and next-year opening remain unavailable. Thus a useful partial history
survives without being presented as a complete year.

The retained adapter has an explicit, narrower role. It recomputes the other
corporate events from their original state and checked methods, verifies that
their P opening balances agree with the shared book, and supplies their
journals and cash origins. Saved expected journals are comparison evidence,
not instructions to the executor. S's book and the group ownership state
remain in that retained method. Replaying them checks compatibility; it does
not claim a second independent implementation of those same corporate rules.

\paragraph{Tax accumulates within a year and resets at its boundary.}
The settlement records its taxable-minus-book adjustment of $-24$ with its
event identity and date. At either tax close, let $B_y(t)$ be P's cumulative
pretax book income from the start of year $y$ through cutoff $t$, let
$A_y(t)$ be the sum of the selected method's dated tax adjustments in that
interval, and let $C_y(t^-)$ be current tax already recognized in the year.
Under the retained constant rate $r=1/4$, the additional current tax is
\begin{equation}
 \Delta C_y(t)=r\{B_y(t)+A_y(t)\}-C_y(t^-).
 \label{eq:return-reference-current-tax}
\end{equation}
The final subtraction prevents an interim tax charge from being paid again.
At the first close, $B_y=56$, $A_y=-24$ and $C_y(t^-)=0$, giving 8. At the
annual close, $B_y=66$ and the same cumulative adjustment remains $-24$;
subtracting the existing 8 gives the remaining 2.5. Negative additional tax
would require a refund or other tax method that this finite reference does
not supply. The opening deferred balances reverse separately, producing
deferred expense 6.

The annual legal statements reconcile profit 49.5, cash 491.5 and equity
671.5. Their operating cash of 101.5 agrees with both the gross payments and
the indirect bridge in equation~\eqref{eq:module-return-cash}. In particular,
the opening interest payment remains a cash adjustment even when current
finance expense is zero. Legal receipts from S stay in P's cash flow; they
are eliminated only by the separate group method.

At year end, the reference transfers profit to retained earnings and clears
the income and period cash accumulators. It retains the source context,
settled return component, disposed-lot identity and lifetime event history.
The old tax adjustment remains available to explain 2025, but its date is
outside $A_{2026}(t)$. In a deliberately inactive following year, equation~
\eqref{eq:return-reference-current-tax} therefore gives zero additional tax
and the deferred reversal is not repeated. This checks annual memory; it is
not a prediction of an inactive business.

The complete reference compares every P account and gross cash origin after
all thirteen dated transitions, then checks the annual statements and carry.
It also preserves the existing six company-authoring paths. The remaining
language task is a composed-company declaration that makes each selected
method, state allocation and unresolved dependency explicit. General changes
in return estimates, partial lot sales, other tax treatments and production
execution still need their own design and conformance cases.

The matched authoring and execution records are under
\path{data/accounting-master/}:
\begin{quote}\small
\path{company-return-schema-20261023}\\
\path{company-return-execution-20261024}
\end{quote}
They retain
the original source comparisons and failure cases. The resulting methods are
conditional finite references; independent implementation, domain review and
human acceptance of the governing design remain pending.

